\section{Conclusion}\label{sec:conclusion}

A CuTe layout is a map that is additive wherever mixed-radix addition does not carry, and the
layout algebra is the calculus of what happens when it does. Seen this way, coalescing deletes
jumps that are zero, composition evaluates one layout at the strides of another and succeeds when
the carries are invisible, CuTe's divisibility conditions single out strides with one digit that
divides its mode, and complements are the classical direct-sum decompositions of an interval. Swizzles fit
the same picture with a second radix: they are layouts for XOR, and XOR differs from integer
addition by the carries of the binary expansion.

\paragraph{In practice.}
For a reader who wants checks rather than theorems, the results amount to four rules.
\begin{enumerate}[leftmargin=1.6em,itemsep=2pt]
\item \emph{Reading a layout.} $\lay SD$ is the function $x\,d_0+\sum_i\floor{x/P_i}\,\jump_i$ with
 $\jump_i=d_i-s_{i-1}d_{i-1}$. On its domain it is determined by its slope and by the total jumps
 at the scales $2\le P<|S|$, and only the non-zero ones matter
 (\cref{prop:floorsum,cor:coalesce}).
\item \emph{Composition.} $A\circ B$ can only be a by-mode refinement $B'$ of $B$, with modes
 $(t'_l,r'_l)$, carrying the strides $\ext A(r'_l)$ at its non-unit modes (\cref{lem:forced};
 \cref{ex:docs} shows the difference between $B$ and $B'$). It is correct if
 $\sum_l(t'_l-1)\,(r'_l\bmod P)<P$ for every scale $P$ of $A$ with a non-zero jump, and this is
 also necessary when the jumps of $A$ are independent, for instance when all non-zero jumps have
 the same sign, as for row-major, column-major and padded matrices (\cref{thm:carryfree}).
 Testing each mode of $B$ on its own, as the implementations do, is not enough when two modes
 land in the same digit of $A$ (\cref{ex:shah}).
\item \emph{Complement and division.} An injective layout with non-negative strides has a
 complement in $\rng M$ exactly when its non-unit modes, sorted by stride, satisfy
 $s_{(i)}d_{(i)}\mid d_{(i+1)}$ and $s_{(k)}d_{(k)}\mid M$, and no other kind of set can tile
 either (\cref{thm:complement}); a layout can be divided by a tiling of its index range exactly
 when its scales with non-zero jump are comparable with the scales of the tiling
 (\cref{cor:divide}).
\item \emph{Swizzles.} If the layout under a swizzle is bit-disjoint---power-of-two mode sizes on
 disjoint bit ranges---and the swizzled layout starts from offset $0$, then the offset bookkeeping
 CuTe uses when slicing it is valid, and a slice is an ordinary layout exactly when no XOR-ed pair
 of bits is split between two of the binary digits that still vary
 (\cref{thm:decay,prop:fields}).
\end{enumerate}

Several questions remain. The carry-free criterion is exact for outer layouts with independent
jumps; for the remaining ones (mixed-sign jumps, torsion codomains) the defect condition of
\cref{thm:exact} is exact but we have no closed form; likewise \cref{thm:decay} is a closed form
for bit-disjoint layouts only. The complement outside the tiling regime,
layouts with symbolic (run-time) extents, and the higher-level operations built on top of
division and product (blocked and raked products, tensor slicing) are not treated here. Finally,
the differential tests suggest concrete changes to implementations. Checking the carry-free
inequality before distributing over the modes of the inner layout would turn the interference
errors into rejections; the single-mode errors of the C++ implementation call for a static stride
test that requires divisibility, as documented; and carry-free compositions with sum strides
could optionally be accepted. For swizzled layouts the corresponding checks are bit-disjointness
of the inner layout when slicing, and the pairwise correspondence of bits of
\cref{prop:transport} when a swizzle is pushed through a layout.

\paragraph{Acknowledgments.}
This work was carried out with substantial assistance from an AI system (Claude, developed by
Anthropic), which, under the author's direction, surveyed the literature, proposed and developed
the framework, wrote the proofs, the Lean formalisation and the test programs, and drafted the
manuscript. Drafts were independently reviewed by a second AI system (Codex, developed by
OpenAI), which found a number of errors that have been corrected and contributed \cref{ex:three},
\cref{ex:assoc} and several of the counterexamples of \cref{sec:swizzle,sec:testing}. The author
takes responsibility for the content. The AI systems are not authors.
